\chapter{Growth Disorders}

\section*{Definition and Overview}
Growth disorders are conditions in which a child's growth is significantly outside the normal range, either too slow (short stature or growth failure) or too fast (excessive growth). Growth is an important indicator of a child's overall health, and deviations from expected patterns can reflect genetic, hormonal, nutritional, or chronic disease factors. Most short or tall children are simply following their genetic potential, but some have an underlying condition that needs assessment. Growth disorders are evaluated by measuring height and weight against growth charts over time, and may involve blood tests, imaging, and referral to a paediatric endocrinologist. Many growth disorders are treatable, and early detection improves outcomes.

\section*{Epidemiology}
Short stature is the most common reason for growth assessment, affecting around 3--5\% of children depending on the definition used. Most short children are healthy and have a normal variant of growth, but around 5--10\% of those referred have an identifiable medical cause. Growth hormone deficiency affects roughly 1 in 4,000 children, and Turner syndrome affects about 1 in 2,500 girls. Excessive growth is less common, with causes such as constitutional tall stature, precocious puberty, and rare syndromes. Growth disorders are diagnosed more often in boys than girls, partly because of referral patterns.

\section*{Aetiology and Pathophysiology}
Growth is regulated by genetics, hormones, and nutrition, and disorders can arise in any of these.

\begin{itemize}
    \item \textbf{Genetic potential:} parental height is the strongest determinant of a child's height
    \item \textbf{Growth hormone:} produced by the pituitary gland, it drives growth in childhood; deficiency causes slow growth, while excess causes gigantism
    \item \textbf{Thyroid hormone:} hypothyroidism slows growth and development
    \item \textbf{Sex hormones:} their rise at puberty drives the growth spurt and eventually closes the growth plates
    \item \textbf{Nutrition:} malnutrition and malabsorption cause growth failure
    \item \textbf{Chronic disease:} conditions such as coeliac disease, inflammatory bowel disease, kidney disease, and asthma can impair growth
    \item \textbf{Chromosomal and genetic conditions:} Turner syndrome, Down syndrome, and many rare syndromes affect growth
    \item \textbf{Constitutional delay:} a common normal variant in which growth and puberty are simply later than average
\end{itemize}

\section*{Clinical Features}
The features depend on the type and cause of the growth disorder.

\begin{itemize}
    \item Height or growth velocity below the 3rd percentile, or crossing down across percentiles over time
    \item Short stature with normal proportions, or with disproportion (short limbs or trunk)
    \item Delayed puberty, which is common in constitutional delay and growth hormone deficiency
    \item Symptoms of hypothyroidism: fatigue, constipation, and cold intolerance
    \item Symptoms of chronic disease: abdominal pain, diarrhoea, or poor appetite
    \item Excessive height or growth velocity above the 97th percentile
    \item Early or delayed puberty
    \item In infants: failure to thrive, with poor weight gain
    \item Features of specific syndromes, such as webbed neck in Turner syndrome
\end{itemize}

\section*{Differential Diagnosis}
The many causes of abnormal growth must be systematically considered.

\begin{itemize}
    \item Constitutional short stature (a small but healthy child)
    \item Constitutional delay of growth and puberty
    \item Familial (genetic) short stature
    \item Growth hormone deficiency and other pituitary disorders
    \item Hypothyroidism
    \item Malnutrition, including from coeliac disease and inflammatory bowel disease
    \item Chronic illness, including kidney disease and asthma
    \item Turner syndrome and other chromosomal conditions
    \item Skeletal disorders, including achondroplasia
    \item Precocious puberty (early puberty) and its effects on final height
    \item Excessive growth: constitutional tall stature, gigantism, Marfan syndrome, and Sotos syndrome
\end{itemize}

\section*{When to Seek Medical Care}
Speak with a GP if a child's growth seems slow, if they are much shorter or taller than peers, if growth velocity has slowed, or if growth problems are accompanied by other symptoms. Early review matters because some conditions are treatable and because growth charts track problems over time.

\textbf{Call 000 (or local emergency services) for:}
\begin{itemize}
    \item Sudden severe headache, vomiting, or visual disturbance in a child with a growth disorder, which may indicate a pituitary problem
    \item Seizures or sudden neurological symptoms
    \item Severe dehydration or acute illness in a child with poor growth
    \item Any emergency related to an underlying condition
\end{itemize}

\section*{Diagnosis and Investigations}
Assessment is methodical, starting with careful measurement and growth chart review.

\begin{itemize}
    \item \textbf{Growth measurement:} height, weight, and head circumference plotted on standard growth charts over time
    \item \textbf{Growth velocity:} the rate of growth is more informative than a single measurement
    \item \textbf{Bone age:} an X-ray of the hand and wrist estimates skeletal maturity and remaining growth potential
    \item \textbf{Blood tests:} thyroid function, growth hormone levels (with stimulation tests), coeliac antibodies, and screens for chronic disease
    \item \textbf{Chromosomal and genetic testing:} for suspected syndromes, including Turner syndrome
    \item \textbf{Imaging:} MRI of the pituitary if growth hormone deficiency is suspected
    \item \textbf{Family and birth history:} parental heights, birth weight, and family patterns of growth
    \item \textbf{Referral:} to a paediatric endocrinologist or growth clinic
\end{itemize}

\section*{Management}
Treatment depends on the cause, and many growth disorders respond well to therapy.

\begin{itemize}
    \item \textbf{Reassurance and monitoring:} most children with short or tall stature need no treatment, just regular measurement
    \item \textbf{Nutritional support:} correcting malnutrition and treating malabsorption
    \item \textbf{Treating underlying disease:} managing coeliac disease, inflammatory bowel disease, kidney disease, and other conditions
    \item \textbf{Thyroid hormone replacement:} for hypothyroidism
    \item \textbf{Growth hormone therapy:} for growth hormone deficiency, Turner syndrome, and other approved indications, given by daily injection
    \item \textbf{Puberty management:} treatment for precocious puberty, and support and reassurance for constitutional delay
    \item \textbf{Monitoring:} regular growth measurements, bone age, and treatment reviews
    \item \textbf{Psychosocial support:} addressing the emotional and social effects of being much shorter or taller than peers
\end{itemize}

\section*{Complications}
\begin{itemize}
    \item Psychosocial effects, including low self-esteem, bullying, and social difficulties
    \item Missed diagnosis of treatable conditions when growth is not monitored
    \item Complications of underlying disease, such as coeliac disease and chronic kidney disease
    \item Side effects of growth hormone therapy, including injection site reactions and fluid retention
    \item Effects of delayed or precocious puberty on final height and psychosocial adjustment
    \item Reduced final height when treatable conditions are diagnosed late
\end{itemize}

\section*{Prognosis and Outcomes}
The outlook depends on the cause. Most short children are healthy and reach a height within their genetic potential. Treatable causes respond well: growth hormone therapy in deficiency adds substantially to final height, and treating hypothyroidism, coeliac disease, and other conditions restores growth. Children diagnosed and treated early achieve much better outcomes than those diagnosed late, which is why growth monitoring at routine health checks is so important. With support, children with growth disorders generally lead full and healthy lives.

\section*{Prevention and Self-Care}
\begin{itemize}
    \item Attend routine health checks so height and weight are plotted regularly
    \item Keep a record of your child's height measurements over time
    \item Ensure good nutrition, including adequate protein, calcium, and vitamin D
    \item Encourage sleep, as growth hormone is largely secreted during deep sleep
    \item Encourage regular physical activity
    \item Seek review if your child is markedly shorter or taller than peers or growth has slowed
    \item Be sensitive to the social and emotional effects of height differences
    \item Follow treatment plans and attend growth clinic reviews
    \item Avoid unproven "growth supplements" and seek professional advice instead
\end{itemize}

\section*{Sources and Further Reading}
The following reputable sources provide further information for healthcare students and patients seeking reliable, current advice about growth disorders.

\begin{itemize}
    \item Australasian Paediatric Endocrine Group. \textit{Growth disorders information.} https://www.apeg.org.au/
    \item Healthdirect Australia. \textit{Growth and development in children.} https://www.healthdirect.gov.au/
    \item Royal Children's Hospital Melbourne. \textit{Growth problems in children.} https://www.rch.org.au/
    \item Turner Syndrome Association of Australia. \textit{Information about Turner syndrome.} https://turnersyndrome.org.au/
    \item NHS. \textit{Growth disorders in children.} https://www.nhs.uk/
    \item American Academy of Pediatrics. \textit{HealthyChildren: growth and development.} https://www.healthychildren.org/
\end{itemize}
